\documentclass[12pt,oneside]{book}

% =====================================================
% METADATOS
% =====================================================
\newcommand{\universidad}{Universidad de Buenos Aires}
\newcommand{\facultad}{Facultad de Derecho}
\newcommand{\carrera}{}
\newcommand{\materia}{Derecho Procesal Civil}
\newcommand{\titulo}{Apuntes de Derecho Procesal Civil}
\newcommand{\autor}{Isaac Edgar Camacho Ocampo}
\newcommand{\anio}{2024}

% =====================================================
% PAQUETES
% =====================================================
\usepackage[utf8]{inputenc}
\usepackage[T1]{fontenc}
\usepackage[spanish,es-nodecimaldot]{babel}

\usepackage[
    top=2.5cm,
    bottom=2.5cm,
    left=3cm,
    right=2.5cm,
    headheight=15pt
]{geometry}

\usepackage{amsmath}
\usepackage{amssymb}
\usepackage{mathtools}

\usepackage{graphicx}
\usepackage{float}
\usepackage{caption}
\usepackage{subcaption}

\usepackage{booktabs}
\usepackage{array}
\usepackage{longtable}
\usepackage{tabularx}

\usepackage{enumitem}

\usepackage{xcolor}
\usepackage[most]{tcolorbox}

\usepackage{fancyhdr}
\usepackage{ragged2e}

\graphicspath{
    {./img/}
    {./graficos/}
    {./figuras/}
}

\usepackage[
    colorlinks=true,
    linkcolor=blue!50!black,
    citecolor=green!40!black,
    urlcolor=blue!60!black,
    pdftitle={\titulo},
    pdfauthor={\autor},
    pdfsubject={Apuntes universitarios},
    bookmarks=true
]{hyperref}

% =====================================================
% CONFIGURACIÓN GENERAL
% =====================================================
\setcounter{tocdepth}{2}

\setlength{\parindent}{0pt}
\setlength{\parskip}{0.5em}

\setlist[itemize]{
    topsep=0.3em,
    itemsep=0.2em
}

\setlist[enumerate]{
    topsep=0.3em,
    itemsep=0.2em
}

\clubpenalty=10000
\widowpenalty=10000

% =====================================================
% ENCABEZADOS Y PIES
% =====================================================
\pagestyle{fancy}
\fancyhf{}

\fancyhead[L]{\nouppercase{\leftmark}}
\fancyhead[R]{\materia}

\fancyfoot[C]{\thepage}

\renewcommand{\headrulewidth}{0.4pt}
\renewcommand{\footrulewidth}{0pt}

\fancypagestyle{plain}{
    \fancyhf{}
    \fancyfoot[C]{\thepage}
    \renewcommand{\headrulewidth}{0pt}
}

% =====================================================
% COMANDOS MATEMÁTICOS
% =====================================================
\newcommand{\abs}[1]{\left|#1\right|}
\newcommand{\norm}[1]{\left\lVert#1\right\rVert}
\newcommand{\vect}[1]{\mathbf{#1}}
\newcommand{\deriv}[2]{\frac{d#1}{d#2}}
\newcommand{\pderiv}[2]{\frac{\partial#1}{\partial#2}}

% =====================================================
% CAJAS ACADÉMICAS
% =====================================================
\newtcolorbox{definition}[1]{
    enhanced,
    breakable,
    title={Definición: #1},
    fonttitle=\bfseries,
    colback=blue!3,
    colframe=blue!50!black,
    boxrule=0.8pt,
    arc=2mm
}

\newtcolorbox{important}{
    enhanced,
    breakable,
    title={Importante},
    fonttitle=\bfseries,
    colback=yellow!8,
    colframe=orange!70!black,
    boxrule=0.8pt,
    arc=2mm
}

\newtcolorbox{example}[1]{
    enhanced,
    breakable,
    title={Ejemplo: #1},
    fonttitle=\bfseries,
    colback=green!4,
    colframe=green!40!black,
    boxrule=0.8pt,
    arc=2mm
}

\newtcolorbox{formula}[1]{
    enhanced,
    breakable,
    title={Fórmula / Regla: #1},
    fonttitle=\bfseries,
    colback=purple!4,
    colframe=purple!50!black,
    boxrule=0.8pt,
    arc=2mm
}

\newtcolorbox{exercise}[1]{
    enhanced,
    breakable,
    title={Ejercicio: #1},
    fonttitle=\bfseries,
    colback=gray!5,
    colframe=gray!60!black,
    boxrule=0.8pt,
    arc=2mm
}

\newtcolorbox{note}{
    enhanced,
    breakable,
    title={Nota},
    fonttitle=\bfseries,
    colback=white,
    colframe=gray!50,
    boxrule=0.6pt,
    arc=2mm
}

% =====================================================
% DOCUMENTO
% =====================================================
\begin{document}

% ---------- PORTADA ----------
\begin{titlepage}

\thispagestyle{empty}

\begin{center}

\vspace*{1cm}

{\large \textsc{\universidad}}

\vspace{0.2cm}

{\large \textsc{\facultad}}

\vspace{3cm}

{\Huge \textbf{\materia}}

\vspace{1cm}

{\LARGE \titulo}

\vspace{0.8cm}

\textit{Comprender el proceso es el primer paso para dominar la práctica del derecho.}

\vspace{1.2cm}

\rule{0.70\textwidth}{0.5pt}

\vspace{0.5cm}

{\large Apuntes de estudio}

\vspace{0.5cm}

\rule{0.70\textwidth}{0.5pt}

\vfill

{\large \autor}

\vspace{0.3cm}

{\large \anio}

\end{center}

\vfill

\begin{flushleft}
\textit{Este es un modesto aporte a los estudiantes de ``Derecho Procesal Civil'' no pretende sustituir el material oficial de la materia.}
\end{flushleft}

\end{titlepage}

% ---------- ÍNDICE ----------
\tableofcontents

% =====================================================
% CAPÍTULO ÚNICO
% =====================================================
\chapter[Litis, audiencia preliminar y prueba]{El proceso civil: litis, audiencia preliminar y sistema probatorio}

\section{La litis, la demanda y la sentencia}

\subsection{El artículo 139 del Código Procesal Civil y Comercial}

Cuando se traba la litis mediante la demanda y su contestación, el juez debe analizar la cuestión planteada. Si esta puede resolverse aplicando exclusivamente el derecho, sin que existan hechos controvertidos, el juez la declara de puro derecho. Si, en cambio, existen hechos en disputa que resulten conducentes, el proceso se abre a prueba y se cita a audiencia preliminar.

\begin{formula}{Artículo 139 -- Código Procesal Civil y Comercial}
Trabada la litis mediante la demanda y su contestación, el juez analiza la cuestión planteada:
\begin{itemize}
    \item Si puede resolverse aplicando únicamente el derecho, sin hechos controvertidos, se declara la cuestión como de puro derecho.
    \item Si existen hechos en disputa y conducentes, el proceso se abre a prueba y se cita a audiencia preliminar.
\end{itemize}
\end{formula}

\begin{definition}{Cuestión de puro derecho}
Es aquella en la que los hechos no importan o no están en discusión. Por ejemplo, si se discute si una ley es constitucional, se trata de una cuestión de puro derecho: no es necesario probar cómo ocurrió algo en la realidad, sino interpretar la Constitución. Por el contrario, cuando se discute cómo ocurrieron los hechos, existen hechos controvertidos que requieren prueba.
\end{definition}

\section{La audiencia preliminar (audiencia de 360 grados)}

La audiencia preliminar es obligatoria y constituye el corazón del proceso civil moderno. En ella ocurren principalmente tres cosas: la conciliación, la determinación del objeto de la litis y la fijación de la prueba.

\begin{table}[H]
\centering
\begin{tabularx}{\textwidth}{
    >{\raggedright\arraybackslash}p{0.28\textwidth}
    >{\raggedright\arraybackslash}X
}
\toprule
\textbf{Etapa} & \textbf{Descripción} \\
\midrule
Conciliación & El juez llama a conciliar. Generalmente no hay acuerdo, pero el juez debe formular la pregunta. \\
Objeto de la litis & El juez identifica los hechos controvertidos y conducentes, y pregunta la versión de cada parte. \\
Fijación de la prueba & Se ofrecen los medios de prueba. El juez va depurando la prueba, admitiéndola o desestimándola según su conducencia. \\
\bottomrule
\end{tabularx}
\caption{Estructura y funciones de la audiencia preliminar}
\end{table}

El objetivo de esta audiencia es que el juez tenga mayor contacto directo con el caso y con las partes: de este modo conoce los hechos de primera mano y se compenetra con el tema. Además, si se realiza correctamente, reduce significativamente la duración del proceso, porque se descartan pruebas impertinentes y redundantes.

\section{Legitimación activa y pasiva}

\begin{definition}{Legitimación activa}
Es la capacidad jurídica de ser demandante. No cualquiera puede demandar a cualquiera por cualquier cosa: es necesario contar con titularidad del derecho, perjuicio actual e interés legítimo.
\end{definition}

\begin{definition}{Legitimación pasiva}
Es la capacidad de ser demandado. Si se demanda a quien no puede responder por la pretensión, la acción es rechazable por falta de legitimación pasiva.
\end{definition}

\section{Caso práctico: Consejo Federal Pesquero}

El Consejo Federal Pesquero es una entidad estatal que otorga licencias de pesca. El Congreso dictó la Ley 24.922 (1998), que crea este consejo y establece que debe otorgar cuotas a las empresas mediante licitación o un mecanismo que garantice la igualdad.

\begin{formula}{Ley 24.922 -- Ley Federal de Pesca}
Crea el Consejo Federal Pesquero, entidad estatal encargada de otorgar licencias de pesca. Establece que las cuotas deben otorgarse a las empresas mediante licitación o un mecanismo que garantice la igualdad.

% TODO: contenido incompleto o ambiguo en las notas originales (el apunte no precisa el artículo específico del decreto reglamentario invocado).
El decreto reglamentario permite al Consejo otorgar cuotas por resolución administrativa. En la práctica, el Consejo otorga cuotas a empresas históricamente favorecidas.
\end{formula}

Una empresa que cumple los requisitos pero ya no accede a la cuota de merluza negra demanda por inconstitucionalidad.

\subsection{El problema constitucional}

Los recursos de pesca son bienes de dominio público. Su explotación debe garantizar igualdad, no discriminación y transparencia. El mecanismo actual viola estos principios porque favorece históricamente a las mismas empresas.

\section{Medios de prueba}

\begin{definition}{Prueba}
Es el conjunto de medios, instrumentos y actuaciones mediante los cuales las partes llevan convicción de sus pretensiones ante el juez.
\end{definition}

\begin{definition}{Fuente de prueba vs. medio de prueba}
Las fuentes de prueba existen en la realidad (un documento, un testigo, un perito). Los medios de prueba son los mecanismos legales para introducir esas fuentes en el proceso. Puede existir una fuente sin que exista un medio legal para ingresarla.
\end{definition}

\begin{table}[H]
\centering
\begin{tabularx}{\textwidth}{
    >{\raggedright\arraybackslash}p{0.32\textwidth}
    >{\raggedright\arraybackslash}X
}
\toprule
\textbf{Medio de prueba} & \textbf{Descripción} \\
\midrule
Documental & Documentos públicos o privados presentados en el expediente. \\
Testimonial & Declaración de testigos que presenciaron los hechos. \\
Informes & Datos técnicos o administrativos provenientes de organismos. \\
Absolución de posiciones & Declaración de la parte en juicio sobre los hechos. \\
Reconocimiento judicial & Inspección directa del lugar por parte del juez. \\
Pericial & Dictamen de peritos expertos en áreas específicas. \\
Medios judiciales & Otros medios que disponga el juez. \\
\bottomrule
\end{tabularx}
\caption{Clasificación de los medios de prueba}
\end{table}

\section{Carga de la prueba y carga dinámica}

\begin{definition}{Carga de la prueba}
Corresponde a quien alega un hecho: el actor prueba sus pretensiones y el demandado prueba sus excepciones.
\end{definition}

\begin{definition}{Carga dinámica de la prueba}
Es el desplazamiento de la carga probatoria hacia quien se encuentra en mejor posición informativa, aun cuando no sea quien alega el hecho.
\end{definition}

\begin{example}{Mala praxis médica}
Un paciente se somete a una cirugía y queda con secuelas. Demanda por mala praxis. El paciente debe probar la lesión (tarea fácil) y que esta fue causada por negligencia médica (tarea difícil). El médico está en mejor posición para probar que actuó correctamente. Por eso, la carga dinámica desplaza hacia el médico: es él quien debe probar que el procedimiento fue correcto.
\end{example}

\begin{important}
La carga dinámica de la prueba no es automática. Debe invocarse activamente en los escritos. Es una herramienta poderosa que muchos abogados dejan pasar.
\end{important}

\section{Sana crítica y apreciación de la prueba}

\begin{definition}{Sana crítica}
Es el sistema mediante el cual el juez valora las pruebas. No es un criterio único ni objetivo, sino un estándar de racionalidad: las conclusiones deben ser lógicas, coherentes y fundadas.
\end{definition}

Si el juez toma una decisión, debe hacerlo paso a paso. Una sentencia correcta debe estructurarse de la siguiente manera: se identifican los hechos controvertidos probados (A, B, C); de la prueba testimonial surge X; de la prueba pericial surge Y; por lo tanto, se concluye Z.

\subsection{La verdad de los hechos frente a la verdad legal}

Entre lo que realmente ocurrió, lo que se prueba y lo que sentencia el juez existen tres cosas distintas. La verdad real es lo que efectivamente pasó. La prueba es un acercamiento a esa verdad. La verdad legal es lo que el juez declara en la sentencia. Con suerte, ambas coinciden; esto depende de la calidad de la prueba y de la habilidad de los abogados.

\section{Prueba pericial: puntos, impugnación y control}

La prueba pericial es fundamental en los procesos complejos. Cuando el juez no cuenta con conocimientos técnicos especializados para determinar los hechos, es necesario que designe un perito experto.

\begin{important}
Al ofrecer prueba pericial se debe ser específico sobre qué debe determinar el perito. Esto se realiza mediante puntos de pericia, que son preguntas concretas. Nunca debe escribirse al final ``y cualquier otro punto de interés'': esto constituye un error grave, viola el debido proceso y hace que la contraparte pierda el derecho a cuestionar los puntos.
\end{important}

\begin{example}{Puntos de pericia en una demanda con pericia psicológica}
Se designa un perito psicólogo para determinar:
\begin{itemize}
    \item Grado o porcentaje de incapacidad psíquica.
    \item Si la incapacidad es permanente o transitoria.
    \item Si requiere tratamiento psicológico.
    \item En qué actividades la persona está impedida de desempeñarse.
    \item Tiempo estimado de tratamiento.
    \item Costo estimado del tratamiento.
\end{itemize}
Cada punto es concreto y mensurable. El perito no tiene opción de decidir qué responder, y la contraparte tiene derecho a impugnar esos puntos si le parecen erróneos.
\end{example}

\section{Reservas de caso federal}

Al redactar una demanda, los abogados suelen incluir la fórmula ``hago reserva de recurrir ante la Corte Suprema''. Esta fórmula es formalmente correcta, pero muchas veces resulta inútil.

\begin{important}
La reserva de caso federal debe ser puntual y específica. No basta con decir ``reservo caso federal'': es necesario identificar cuál es la cuestión constitucional en juego, cuál artículo constitucional se invoca y por qué. Aproximadamente el 99\% de las demandas incluyen esta cláusula, pero la mayoría no sirve porque no reservan nada específico. Los jueces no objetan nada al momento de presentarla, pero la reserva resulta inútil al recurrir ante la Corte si no fue formulada correctamente.
\end{important}

\section{Síntesis: responsabilidad y estrategia del abogado}

La respuesta a si es posible incluir en una demanda un pedido ``no usual'', como una prueba extraordinaria, depende de cómo se lo plantee. Son los abogados quienes impulsan los cambios en la práctica judicial: los jueces no modifican la práctica por iniciativa propia. Es mediante el planteo constante de una necesidad que la práctica termina adaptándose.

La carga dinámica de la prueba es un ejemplo de este fenómeno: si se invoca constantemente en cada caso donde resulte aplicable, argumentando por qué esa parte está en mejor posición de probar, eventualmente termina consolidándose como jurisprudencia. Esto es especialmente cierto en materia laboral, donde el código de procedimiento civil se aplica de manera supletoria y el empleado se encuentra habitualmente en desventaja informativa frente al empleador.

El derecho procesal civil es el sistema mediante el cual la justicia resuelve los conflictos entre ciudadanos. Su lógica es secuencial: la demanda (donde se plantea la litis), la audiencia preliminar (donde se organiza la prueba), la producción de pruebas (documental, testimonial, pericial), la apreciación según la sana crítica (racionalidad) y la sentencia (donde el juez declara la verdad legal). Cada etapa tiene reglas que protegen garantías procesales. Como abogado, el rol consiste en conocer profundamente estas reglas, usarlas estratégicamente, invocar activamente institutos como la carga dinámica cuando corresponda y fundamentar siempre la argumentación para poder recurrir en caso de arbitrariedad.

% =====================================================
% CORNELL
% =====================================================
\section*{Cuadro Cornell: síntesis y revisión}
\addcontentsline{toc}{section}{Cuadro Cornell: síntesis y revisión}

\begin{longtable}{
    >{\raggedright\arraybackslash}p{0.30\textwidth}
    >{\raggedright\arraybackslash}p{0.60\textwidth}
}
\toprule
\textbf{Preguntas / palabras clave} & \textbf{Notas / respuestas} \\
\midrule
\endfirsthead

\toprule
\textbf{Preguntas / palabras clave} & \textbf{Notas / respuestas} \\
\midrule
\endhead

\textbf{¿Qué es la litis?} &
La litis es la controversia entre dos partes. Existe desde que hay demanda y contestación. El juez determina si se resuelve solo por derecho o si hay hechos controvertidos que requieren prueba. \\

\textbf{¿Cuál es el objetivo de la audiencia preliminar?} &
Conciliación entre las partes, determinación del objeto del litigio, identificación de hechos controvertidos y fijación de la prueba. Todo esto organiza el proceso y lo vuelve más eficiente. \\

\textbf{¿Quién tiene la carga de la prueba?} &
Quien alega un hecho tiene la carga de probarlo: el actor prueba sus pretensiones y el demandado sus excepciones. Existe la carga dinámica: quien está en mejor posición de probar puede ser obligado a hacerlo. \\

\textbf{¿Qué es la sana crítica?} &
Es el estándar de racionalidad mediante el cual el juez valora las pruebas. Exige que se fundamente paso a paso cómo se llegó a las conclusiones. Una sentencia arbitraria, sin fundamento lógico, es recurrible. \\

\textbf{¿Cómo se construye una demanda por inconstitucionalidad?} &
Identificando el demandado correcto (quien aplica la norma inconstitucional), la legitimación activa (perjuicio actual), la norma inconstitucional y el fundamento (incompatibilidad constitucional). \\

\textbf{¿Qué error grave existe con los puntos de pericia?} &
No ser específico: dejar puntos abiertos como ``y cualquier otro que el perito considere''. Esto viola el debido proceso. Los puntos deben ser concretos y mensurables. \\

\textbf{¿Cuándo se aplica la carga dinámica de la prueba?} &
Cuando una parte tiene acceso informativo que la otra no posee. Ejemplos: mala praxis (el médico está en mejor posición) y derecho laboral (el empleador tiene los datos). Debe invocarse activamente. \\

\textbf{¿Cuál es la diferencia entre fuente y medio de prueba?} &
La fuente es el elemento que existe en la realidad (documento, testigo, experticia). El medio es el mecanismo legal para introducirla en el proceso. Puede existir una fuente sin un medio legal para incorporarla. \\

\textbf{¿Qué diferencia hay entre verdad real, prueba y verdad legal?} &
La verdad real es lo que ocurrió. La prueba es un acercamiento a esa verdad. La verdad legal es lo que el juez declara en la sentencia. Idealmente, ambas coinciden. \\

\textbf{¿Cómo debe estructurarse una sentencia correcta?} &
Identificando los hechos controvertidos probados (A, B, C); señalando que de la prueba testimonial surge X y de la prueba pericial surge Y; concluyendo que, por lo tanto, se consideran probados P, Q, R, y que, aplicando el artículo correspondiente, corresponde la conclusión Z. Todo debe estar conectado lógicamente. \\

\midrule

\multicolumn{2}{p{0.92\textwidth}}{
\textbf{Resumen:} El proceso civil se organiza secuencialmente a partir de la litis planteada en la demanda y su contestación. La audiencia preliminar concilia, fija el objeto del litigio y organiza la prueba. Rigen la carga de la prueba y, excepcionalmente, la carga dinámica cuando una parte tiene mejor acceso a la información. La prueba pericial exige puntos concretos y mensurables. El juez aprecia la prueba conforme a la sana crítica, un estándar de racionalidad que exige fundamentación lógica paso a paso, y declara en la sentencia una verdad legal que debería acercarse a la verdad real. La reserva de caso federal solo resulta útil si identifica de manera específica la cuestión y el artículo constitucional en juego.
} \\

\bottomrule
\end{longtable}

\end{document}
